\chapter{Le Monde Adulte : Le Théâtre des Masques}

\textit{Devenir adulte, c'est souvent troquer ses rêves de liberté contre une armure de convenances. Mais quand cette armure devient trop lourde, quand elle nous empêche de respirer et de dire qui nous sommes, elle finit par nous étouffer.}

\section{L'Entreprise : La Machine à Inhiber}
Le monde du travail est une arène où l'affirmation de soi est à la fois exigée et réprimée. On demande du leadership, mais on sanctionne souvent l'impertinence. Pour l'adulte inhibé, l'open-space est une jungle de signaux qu'il doit décoder sans cesse. Le modèle de Karasek nous l'explique : le stress n'est pas dû à la charge de travail seule, mais au manque de contrôle sur celle-ci. En ne sachant pas dire "non" à une tâche supplémentaire, le salarié perd son autonomie et entre dans le tunnel du burnout.

\begin{NoteClinique}
\textbf{Le cas de Franck, 45 ans.} \\
Franck est un "pilier" dans son agence. Il ne refuse rien. On dit de lui qu'il est "fiable". Mais Franck est vide. Chaque fois qu'il accepte un dossier de trop par peur de décevoir son manager, il sent une petite partie de lui s'effondrer. Physiquement, il vit dans un état de fatigue que même deux semaines de vacances ne soignent plus. Franck ne souffre pas de surmenage, il souffre d'un manque d'existence. Son corps a fini par dire "stop" là où son esprit n'a pas osé le faire.
\end{NoteClinique}

\section{Genre et Socialisation : Le Poids des Attendus}
Nous ne sommes pas égaux devant l'affirmation de soi. Les siècles de socialisation pèsent sur nos épaules. On a souvent appris aux femmes que la douceur était une vertu suprême, assimilant l'affirmation de soi à une forme d'agressivité "peu féminine". À l'inverse, on a imposé aux hommes un masque de force imperturbable, où avouer un besoin ou une limite est perçu comme une faiblesse. Ces carcans invisibles créent des solitudes immenses au sein même des couples et des familles.

\section{L'Asymétrie du Silence dans le Couple}
Dans l'intimité, la non-affirmation prend des formes plus subtiles mais tout aussi dévastatrices. On ne dit pas son besoin, on espère que l'autre le devinera. Et quand l'autre ne devine pas, le ressentiment s'installe. Ce "silence bruyant" érode la complicité. S'affirmer dans le couple, ce n'est pas imposer sa volonté, c'est offrir à l'autre la vérité de ce que nous sommes. Sans cette vérité, le lien ne repose que sur des masques qui finissent par se lasser l'un de l'autre.
